% The same protocol on a bus with four joints, and the load arithmetic.
\begin{tikzpicture}[font=\sffamily\small, x=1cm, y=1cm]

% ---------------- the trunk
\draw[busline] (-6.2,0) -- (6.2,0);
\node[busterm] at (-6.2,0) {120};
\node[busterm] at (6.2,0) {120};

% ---------------- five participants
\node[busmaster, text width=26mm] (mm) at (-4.6,1.9) {the master\\{\scriptsize node 0}};
\draw[busstub] (-4.6,0) -- (mm.south);
\foreach \i/\x/\n in {1/-1.6/1, 2/0.6/2, 3/2.8/3, 4/5.0/4}{
  \node[busnode, text width=20mm] (j\i) at (\x,1.9) {joint \n\\{\scriptsize node \n}};
  \draw[busstub] (\x,0) -- (j\i.south);
  \node[font=\sffamily\tiny, anchor=north] at (\x,-0.15) {0x10\n, 0x18\n};
}
\node[font=\sffamily\tiny, anchor=north] at (-4.6,-0.15) {sends 0x10n to each};

\node[chlabel, anchor=west] at (-6.2,2.9)
  {each node's identifiers are its function codes offset by its node number: nothing is configured centrally};

% ---------------- the arithmetic
\node[regcell, minimum width=72mm, text width=70mm, anchor=north west,
      fill=white, align=left, minimum height=30mm] at (-6.4,-1.2)
  {\ttfamily\tiny
   state frame\\
   \hspace*{1em} 67 arb bits @ 500k + 240 data bits @ 2M = 254 us\\
   command frame\\
   \hspace*{1em} 67 arb bits @ 500k + 176 data bits @ 2M = 222 us\\[3pt]
   four joints, both directions, at 1 kHz:\\
   \hspace*{1em} 4 * (254 + 222) us = 1.90 ms per 1.00 ms period\\[3pt]
   \normalfont\sffamily\scriptsize\color{red!60!black}\bfseries
   190 per cent. It does not fit.};

\node[umlnote, text width=54mm, anchor=north west] at (1.2,-1.2)
  {\textbf{Three ways out, in order of preference.} Command at 250 hertz with
   interpolation on the node brings it to 91 per cent. Raising the arbitration
   rate to one megabit brings it to 61 per cent. Both together give 38 per cent,
   which is the design this volume would build for a real arm.\\[3pt]
   The arbitration phase is the expensive part, because it runs at the slow
   rate on every frame regardless of payload. That is why raising the
   arbitration rate helps more than shrinking the layout.};

\node[chlabel, anchor=north west, text width=112mm] at (-6.4,-6.4)
  {This bench has one joint and a master, so the four-joint figure is
   arithmetic rather than a measurement, and the chapter says so. What the
   arithmetic is good for is deciding the design before the second joint
   exists, which is exactly when it is cheap to change.};
\end{tikzpicture}
